%==================================================================
% Ini adalah bab 3
% Silahkan edit sesuai kebutuhan, baik menambah atau mengurangi \section, \subsection
%==================================================================

\chapter[METODE PENELITIAN]{\\ METODE PENELITIAN}

Minimal 8 halaman. Silakan gunakan template yang sesuai dengan tema laporan anda.\\

%==================================================================
% Tema Networking/Jaringan
%==================================================================
\textbf{Tema Networking/Jaringan}

\setcounter{section}{0} % Silakan hapus jika sudah mulai menulis
\setcounter{subsection}{0} % Silakan hapus jika sudah mulai menulis

\section{Pembahasan}

\subsection{Jenis Jaringan}
\begin{subsectioncontent}
    \hspace{\parindent}Teori singkat tentang jenis jaringan yang digunakan pada sistem yang dibahas (Peer to Peer, client Server, Network Based, atau bentuk khusus lain). Berikan penjelasan tentang kelebihan dan kekurangannya untuk jenis yang digunakan pada tempat KKP.
\end{subsectioncontent}

\subsection{Topologi}
\begin{subsectioncontent}
    \hspace{\parindent}Teori singkat tentang topologi jaringan yang digunakan pada sistem yang dibahas (Topologi Star, Topologi Bus, Topologi Ring atau bentuk khusus lainnya). Berikan penjelasan tentang kelebihan dan kekurangannya untuk topologi yang digunakan pada tempat KKP.
\end{subsectioncontent}

\subsection{Gambar Jaringan Logical}
\begin{subsectioncontent}
    \hspace{\parindent}Gambarkan kondisi jaringan saat ini dengan lengkap beserta penjelasan/terkait. (gunakan visio).
\end{subsectioncontent}

\section{Gambar Jaringan Fisik}
\begin{sectioncontent}
    \hspace{\parindent}Gambarkan posisi perangkat jaringan yang ada pada gambar logical, pada tempat KKP. Misalnya, server pada lantai II, Finger Print pada lantai III, dst. (Gunakan visio)
\end{sectioncontent}

\section{Pengalamantan}
\begin{sectioncontent}
    \hspace{\parindent}Jelaskan jenis IP (versi 4 /6), kelas pengalamatan, dan pembagian pengalamatan dan subnetting dari jaringan yang ada. (jika organisasi menutup informasi ini, maka subtitusi angka Ipnya, namun gunakan konfigurasi pengalamatan yang sama)
\end{sectioncontent}

\section{Aplikasi Jaringan}
\begin{sectioncontent}
    \hspace{\parindent}Berikan daftar aplikasi dan layanan yang telah disebutkan pada BAB III lalu jelaskan secara detil konfigurasinya.
\end{sectioncontent}

%==================================================================
% Tema UI/UX
%==================================================================
\textbf{Tema UI/UX}

\setcounter{section}{0} % Silakan hapus jika sudah mulai menulis
\setcounter{subsection}{0} % Silakan hapus jika sudah mulai menulis

\section{Analisis kebutuhan pengguna}

\subsection{Spesifikasi Konteks Penggunaan}
\begin{subsectioncontent}
    \hspace{\parindent}Tahapan ini merupakan mengidentifikasi orang yang akan menggunakan produk. Kegiatan ini akan menjelaskan untuk apa dan dalam kondisi seperti apa mereka akan menggunakan produk.
\end{subsectioncontent}

\subsection{Spesifikasi Kebutuhan Pengguna dan Organisasi}
\begin{subsectioncontent}
    \hspace{\parindent}Tahapan ini akan melakukan identifikasi kebutuhan pengguna. Kebutuhan pengguna yang akan teridentifikasi merupakan bagian dari tahapan hasil wawancara.
\end{subsectioncontent}

\subsection{Desain Proses}
\begin{subsectioncontent}
    \hspace{\parindent}Tahapan ini akan membuat rancangan analisis menggunakan diagram UML (Unified Modeling Language) seperti Use Case Diagram dan Activity Diagram.
\end{subsectioncontent}

\section{Perancangan}

\subsection{Rancangan Urutan Komponen Dialog}
\begin{subsectioncontent}
    \hspace{\parindent}Komponen dialog atau ragam dialog merupakan proses pertukaran komunikasi antara satu atau lebih dalam mengambil kaidah informasi dalam user dan sistem komputer.
\end{subsectioncontent}

\subsection{Wireframe}
\begin{subsectioncontent}
    \hspace{\parindent}Pada tahap ini dilakukan pembuatan desain prototype dalam bentuk konsep interface layout yang akan diterapkan pada proses prototype. Untuk pembuatan wireframe dapat menggunakan tool seperti pencil project, figma, dsb.
\end{subsectioncontent}

\subsection{Prototype}
\begin{subsectioncontent}
    \hspace{\parindent}Pada tahap ini dari segi interface nya prototype akan dibuat jauh lebih detail lagi dibandingkan dengan metode wireframe. Dari segi visual dan kontennya sudah lebih berwarna dan hampir menyamai dengan final produk. Pada prototype ini juga disertai dengan transisi dan animasi antar menu serta fitur yang lebih interaktif dan clickable sehingga pengguna bisa menguji dan merasakan perancangan prototype yang telah selesai dibuat.
\end{subsectioncontent}

\section{Pengujian}
\begin{sectioncontent}
    \hspace{\parindent}Pada tahap ini dilakukan pengujian pada Prototype yang telah dibuat sebelumnya dan dilakukan secara mandiri. Hal ini dilakukan untuk memastikan MVP yang telah dibuat sudah sesuai dan dapat berjalan dengan baik sebelum pengujian kepada pengguna.
\end{sectioncontent}

%==================================================================
% Tema Game
%==================================================================
\textbf{Tema Game}

\setcounter{section}{0} % Silakan hapus jika sudah mulai menulis
\setcounter{subsection}{0} % Silakan hapus jika sudah mulai menulis

\section{Perancangan}

\subsection{Latar Belakang Cerita (background story)}
\begin{subsectioncontent}
    \hspace{\parindent}Sub-bab ini berisi cerita singkat yang melatar belakangi game Anda. Biasanya ada pada game ber-genre RPG, Adventure, dan Arcade terdapat alur cerita yang membentuk game tersebut
\end{subsectioncontent}

\subsection{Rincian Game}
\begin{subsectioncontent}
    \hspace{\parindent}Sub-bab ini berisi detail dari game yang Anda buat. Contoh: bila Anda membuat game RPG, “bertema peperangan antar tokoh pada dua kerajaan yang memperebutkan sesuatu” maka subbab ini akan berisi rincian senjata, baju zirah, karakter tokoh beserta statistiknya, kemampuan karakter, dan lain-lain
\end{subsectioncontent}

\section{Arsitektur Game}
Sub-bab ini berisi Use Case Diagram + Class Diagram + Activity/State Diagram + Sequence Diagram (bilamana menggunakan UML) atau Flowchart game.

\subsection{Storyboard}
\begin{subsectioncontent}
    \hspace{\parindent}Sub-bab ini berisi desain rancangan rangkaian cerita dan layout dari game Anda
\end{subsectioncontent}

\subsection{Creative Strategy}
\begin{subsectioncontent}
    \hspace{\parindent}Sub-bab ini menjelaskan design (visual) atau layout yang diimplementasikan pada game Anda
\end{subsectioncontent}

\section{Pengujian}
\begin{sectioncontent}
    \hspace{\parindent}Sub  ini akan menjabarkan hasil uji coba game tersebut dalam bentuk laporan pengujian, dan pembahasan dari tiap class/method/fungsi utama yang dibuat. Pengujian untuk Game project hanya berupa whitebox testing
\end{sectioncontent}

%==================================================================
% Tema Sistem Informasi
%==================================================================
\textbf{Tema Sistem Informasi}

\setcounter{section}{0} % Silakan hapus jika sudah mulai menulis
\setcounter{subsection}{0} % Silakan hapus jika sudah mulai menulis

\section{Proses Bisnis}
\begin{sectioncontent}
    Menceritakan bagaimana sistem bekerja atau di proses dari sistem yang dibahas
    \begin{enumerate}[leftmargin=0.5cm]
        \item Sistem Akuntansi / prosedur akuntansi
        \item Proses Bisnis beserta dengan Activity Diagram
        \item Use Case Diagram
    \end{enumerate}
\end{sectioncontent}

\section{Interface}
\begin{sectioncontent}
    Menceritakan tentang interface dari sistem yang dibahas dengan penjelasannya
    \begin{enumerate}[leftmargin=0.5cm]
        \item Model Layar / Tampilan (gunakan screenshot bila diijinkan)
        \item Model Masukan (Lampirkan Dokumen Masukan)
        \item Model Keluaran (Lampirkan Dokumen Keluaran)
    \end{enumerate}
\end{sectioncontent}

\section{Rancangan Database}
\begin{sectioncontent}
    Menceritakan tentang struktur database dari sistem yang dibahas
    \begin{enumerate}[leftmargin=0.5cm]
        \item Class Diagram /LRS
        \item Spesifikasi basis data (jika organisasi tidak mengetahui spesifikasi basis data, buatlah reka ulang spesifikasi basis data berdasarkan informasi dokumen masukan keluaran dan model proses)
    \end{enumerate}
\end{sectioncontent}